\documentclass[12pt,a4paper]{article}
\usepackage[utf8]{inputenc}
\usepackage[german]{babel}
\usepackage{amsmath, amssymb, amsthm}
\usepackage{geometry}
\usepackage{csquotes}
\geometry{margin=2.5cm}

\newtheorem{axiom}{Axiom}
\newtheorem{definition}{Definition}
\newtheorem{satz}{Satz}
\newtheorem{theorem}{Theorem}
\newtheorem{korollar}{Korollar}
\newtheorem{bemerkung}{Bemerkung}
\newtheorem{diskussion}{Diskussion}
\renewcommand{\proofname}{Beweis}

\title{Die Theorie der fraktalen Entropie\\
\Large Eine konsistente Thermodynamik ohne Wärmetod\\
\Large Mit einer Diskussion der topologischen Optionen}
\author{Michael Czybor}
\date{März 2026}

\begin{document}

\maketitle

\begin{abstract}
Die Theorie der fraktalen Entropie ist eine fundamentale Erweiterung der 
Weber-de-Broglie-Bohm-Theorie mit fraktalem Raum (WDBT+). Sie basiert auf 
sieben Axiomen, die einen geschlossenen, stationären Kreislauf der 
fraktalen Entropie beschreiben. Die Gravitation erzeugt durch die 
DSTT-Drift ($\dot{\beta}>0$) irreversible Prozesse und einen intrinsischen 
Zeitpfeil, während die Elektrodynamik ($\dot{\beta}=0$) stabile 
Strukturen bewahrt. Die fraktale Raumdimension $D\approx2,71$ verhindert 
ein globales Entropiemaximum und damit den Wärmetod. Energie, Masse und 
Dichte bleiben konstant; die fraktale Entropie zirkuliert zwischen 
Vakuum, Materie und $\Omega$-Feld. Die Theorie wird für beide 
topologischen Möglichkeiten (geschlossenes oder offenes Universum) 
diskutiert. In beiden Fällen expandiert der Raum nicht; lediglich 
Strukturen und Bahnen dehnen sich durch die DSTT-Drift aus, kompensiert 
durch kontinuierliche Materieentstehung.
\end{abstract}

\tableofcontents

\section{Einleitung}

Die Weber-de-Broglie-Bohm-Theorie mit fraktalem Raum (WDBT+) beschreibt 
die Gravitation als direkte, instantane Teilchen-Teilchen-Wechselwirkung 
nach dem Weber-Gesetz, gesteuert durch das Bohmsche Quantenpotential, in 
einem Raum mit fraktaler Dimension $D\approx2,71$. Die vorliegende 
Theorie der fraktalen Entropie ergänzt dieses Bild durch eine 
konsistente thermodynamische Fundierung. Ihr zentrales Ergebnis ist, 
dass trotz permanenter irreversibler Prozesse und eines intrinsischen 
Zeitpfeils kein Wärmetod eintritt – sowohl in einem topologisch 
geschlossenen als auch in einem topologisch offenen Universum.

\section{Die Axiome der Theorie}

\subsection{Axiom 1: Fraktale Raumstruktur}

\begin{axiom}[Fraktale Raumstruktur]
Der fundamentale Raum ist eine fraktale Menge $\mathcal{M}$ mit Hausdorff-Dimension $D$, wobei
\begin{equation}
\boxed{D = \frac{\ln 20}{\ln(2 + \phi)} \approx 2,71}, \qquad \phi = \frac{1+\sqrt{5}}{2}
\end{equation}
Die fundamentale Längenskala ist
\begin{equation}
\boxed{l_0 = \left( \frac{V_{\text{Dodekaeder}}}{V_{\text{Einheitskugel}}} \right)^{1/3} \lambda_p \approx 1,8 \times 10^{-15} \, \text{m}}
\end{equation}
mit $\lambda_p$ der Compton-Wellenlänge des Protons. $l_0$ ist die kleinste physikalisch mögliche Distanz – eine natürliche untere, aber keine obere Grenze.
\end{axiom}

\begin{bemerkung}
Die fraktale Dimension $D<3$ ist entscheidend für die gesamte Theorie. Sie bewirkt, dass die fraktale Entropie subextensiv ist und kein globales Maximum besitzt.
\end{bemerkung}

\subsection{Axiom 2: Fraktale Entropie}

\begin{axiom}[Fraktale Entropie]
Für eine Wahrscheinlichkeitsdichte $\rho: \mathcal{M} \to \mathbb{R}^+$ mit $\int_{\mathcal{M}} \rho \, d^D r = 1$ ist die fraktale Entropie definiert als
\begin{equation}
\boxed{S_D[\rho] = -k_B \int_{\mathcal{M}} d^D r \, \rho(\vec{r}) \ln \rho(\vec{r}) \cdot \left( \frac{|\vec{r}|}{l_0} \right)^{D-3}}
\end{equation}
\end{axiom}

\begin{bemerkung}
Für $D=3$ ergibt sich die Boltzmann-Entropie. Für $D<3$ ist $S_D$ subextensiv: Sie wächst langsamer als das Volumen. Dies ist eine direkte Konsequenz der fraktalen Geometrie.
\end{bemerkung}

\subsection{Axiom 3: Globaler Entropieerhaltungssatz}

\begin{axiom}[Globaler Entropieerhaltungssatz]
Die Gesamt-fraktale Entropie des Universums ist konstant:
\begin{equation}
\boxed{\frac{dS_D^{\text{ges}}}{dt} = 0}
\end{equation}
\end{axiom}

\begin{bemerkung}
Die fraktale Entropie zirkuliert zwischen den Phasen (Vakuum, Materie, $\Omega$-Feld), aber ihr Gesamtwert ist zeitlich konstant. Dies ist das fundamentale Stationaritätsaxiom.
\end{bemerkung}

\subsection{Axiom 4: Lokale Entropieproduktion}

\begin{axiom}[Lokale Entropieproduktion]
In jeder einzelnen Phase gilt der zweite Hauptsatz in erweiterter Form:
\begin{equation}
\boxed{\frac{dS_D^{(i)}}{dt} \geq 0}
\end{equation}
Die Entropie wächst lokal, aber durch Zirkulation ist der globale Zuwachs null.
\end{axiom}

\subsection{Axiom 5: Gravitation als Quelle der Irreversibilität}

\begin{axiom}[Gravitation]
Die Weber-Gravitation mit $\beta_{\text{WG}} = 0,5$ führt zu einer monotonen Zunahme des tangentialen Energieanteils:
\begin{equation}
\boxed{\dot{\beta}_{\text{DSTT}} > 0 \quad \text{für alle } t}
\end{equation}
Dies erzeugt eine permanente Auswärtsdrift, irreversible Prozesse und einen intrinsischen Zeitpfeil.
\end{axiom}

\subsection{Axiom 6: Elektrodynamik als Stabilisator}

\begin{axiom}[Elektrodynamik]
Die Weber-Elektrodynamik mit $\beta_{\text{WED}} = 2$ führt zu keiner Änderung der Energieaufteilung:
\begin{equation}
\boxed{\dot{\beta}_{\text{DSTT}} = 0 \quad \text{für alle } t}
\end{equation}
Sie ermöglicht reversible Prozesse, stabile Bahnen (Ellipsen) und bewahrt Strukturen.
\end{axiom}

\subsection{Axiom 7: Beschränktheit aller physikalischen Größen}

\begin{axiom}[Beschränktheit]
Für jede physikalische Größe $X$ (Energie, Masse, Dichte, Temperatur, Volumen, Komplexität) gilt:
\begin{equation}
\boxed{\limsup_{t \to \infty} |X(t)| < \infty}
\end{equation}
Keine physikalische Größe divergiert oder wächst unbegrenzt. Das Universum bleibt in allen messbaren Aspekten stationär.
\end{axiom}

\section{Die Struktur der fraktalen Entropie}

\subsection{Die drei Phasen}

Das Universum besteht aus drei Phasen, zwischen denen die fraktale Entropie zirkuliert:

\begin{definition}[Phasen]
\begin{itemize}
    \item \textbf{Vakuum (Nullpunkt):} $S_D^{\text{vac}}$ – die fraktale Struktur des Raums selbst
    \item \textbf{Materie (kinetisch):} $S_D^{\text{kin}}$ – Bewegung der Teilchen, unterteilt in:
    \begin{itemize}
        \item $S_D^{\text{kin,grav}}$ – durch Gravitation gebundene Materie
        \item $S_D^{\text{kin,em}}$ – durch Elektrodynamik gebundene Materie
    \end{itemize}
    \item \textbf{$\Omega$-Feld:} $S_D^{\Omega}$ – das Gravitationsfeld als Rotationsfeld
\end{itemize}
\end{definition}

Der Erhaltungssatz lautet:
\begin{equation}
S_D^{\text{vac}}(t) + S_D^{\text{kin,grav}}(t) + S_D^{\text{kin,em}}(t) + S_D^{\Omega}(t) = S_D^{\text{ges}} = \text{const}
\end{equation}

\subsection{Die Entropieflüsse}

Aus den Bewegungsgleichungen der WDBT+ folgt:

\begin{definition}[Entropieflussraten]
\begin{align}
\frac{dS_D^{\text{vac}}}{dt} &= -\sigma_{\text{cre}} + \sigma_{\text{back}} \\
\frac{dS_D^{\text{kin,grav}}}{dt} &= +\sigma_{\text{cre}} - \sigma_{\text{drift}} \\
\frac{dS_D^{\text{kin,em}}}{dt} &= 0 \quad \text{(elektromagnetische Prozesse sind reversibel)} \\
\frac{dS_D^{\Omega}}{dt} &= +\sigma_{\text{drift}} - \sigma_{\text{back}}
\end{align}
mit
\begin{itemize}
    \item $\sigma_{\text{cre}} > 0$: Entropiefluss von Vakuum zu gravitativer Materie (Materieentstehung durch Quantenpotential)
    \item $\sigma_{\text{drift}} > 0$: Entropiefluss von gravitativer Materie zum $\Omega$-Feld (DSTT-Drift)
    \item $\sigma_{\text{back}} > 0$: Entropiefluss vom $\Omega$-Feld zum Vakuum (Rückkopplung)
\end{itemize}
\end{definition}

\begin{satz}[Stationärer Zustand]
Im stationären Zustand gilt:
\begin{equation}
\boxed{\sigma_{\text{cre}} = \sigma_{\text{drift}} = \sigma_{\text{back}} \equiv \sigma}
\end{equation}
Damit sind die Entropien jeder Phase konstant:
\begin{equation}
\frac{dS_D^{\text{vac}}}{dt} = \frac{dS_D^{\text{kin,grav}}}{dt} = \frac{dS_D^{\text{kin,em}}}{dt} = \frac{dS_D^{\Omega}}{dt} = 0
\end{equation}
Es zirkuliert nur der Fluss, nicht die Menge.
\end{satz}

\begin{figure}[ht]
\centering
\framebox{\parbox{0.9\textwidth}{\centering
\vspace{0.5cm}
\textbf{Der geschlossene Entropiekreislauf:} \\
Vakuum $S_D^{\text{vac}}$ $\xrightarrow{\sigma_{\text{cre}}}$ Materie (grav.) $S_D^{\text{kin,grav}}$ $\xrightarrow{\sigma_{\text{drift}}}$ $\Omega$-Feld $S_D^{\Omega}$ $\xrightarrow{\sigma_{\text{back}}}$ Vakuum \\
Elektromagnetische Materie $S_D^{\text{kin,em}}$ ist vom Kreislauf entkoppelt (reversibel) \\
Gesamt: $S_D^{\text{ges}} = \text{const}$
\vspace{0.5cm}
}}
\caption{Der geschlossene Kreislauf der fraktalen Entropie.}
\end{figure}

\section{Die duale Dynamik: Gravitation vs. Elektrodynamik}

\subsection{Herleitung aus der erweiterten Weber-Theorie}

Die allgemeine Weber-Kraft lautet:
\begin{equation}
\vec{F} = \frac{Q_1Q_2}{4\pi\epsilon_0 r^2} \left\{ \left[ 1 - \frac{\dot{r}^2}{c^2} + \beta_{\text{Kraft}} \frac{r\ddot{r}}{c^2} \right] \hat{r} + \frac{2(\dot{r} \cdot \vec{v})}{c^2} \vec{v} \right\}
\end{equation}

Für die zeitliche Änderung der Energieaufteilung folgt:
\begin{equation}
\dot{\beta}_{\text{DSTT}} = \frac{2\dot{r}}{r} \cdot \frac{Q_1Q_2}{4\pi\epsilon_0 c^2 r m} \cdot \frac{2 - \beta_{\text{Kraft}}}{\left(1 + \frac{\dot{r}^2}{r^2\dot{\phi}^2}\right)^2} + \mathcal{O}\left(\frac{1}{c^4}\right)
\end{equation}

\begin{satz}[Dualität der Wechselwirkungen]
Der Faktor $(2 - \beta_{\text{Kraft}})$ bestimmt das Verhalten:
\begin{equation}
\boxed{\dot{\beta}_{\text{DSTT}} \begin{cases}
> 0, & \text{falls } \beta_{\text{Kraft}} < 2 \\
= 0, & \text{falls } \beta_{\text{Kraft}} = 2
\end{cases}}
\end{equation}
Daraus folgt:
\begin{itemize}
    \item \textbf{Gravitation:} $\beta_{\text{WG}} = 0,5 \;\Rightarrow\; \dot{\beta}_{\text{DSTT}} > 0$ (Drift, Irreversibilität)
    \item \textbf{Elektrodynamik:} $\beta_{\text{WED}} = 2 \;\Rightarrow\; \dot{\beta}_{\text{DSTT}} = 0$ (keine Drift, Reversibilität)
\end{itemize}
\end{satz}

\begin{table}[ht]
\centering
\caption{Die duale Natur der Wechselwirkungen}
\begin{tabular}{|l|c|c|c|}
\hline
\textbf{Wechselwirkung} & $\beta_{\text{Kraft}}$ & $\dot{\beta}_{\text{DSTT}}$ & \textbf{Rolle} \\ \hline
Gravitation & $0,5$ & $> 0$ & Motor der Zeit, Irreversibilität, Drift \\ \hline
Elektrodynamik & $2$ & $= 0$ & Stabilisator, Reversibilität, Strukturerhalt \\ \hline
\end{tabular}
\end{table}

\begin{bemerkung}
Die Gravitation erzeugt permanente irreversible Flüsse und damit den Zeitpfeil. Die Elektrodynamik ermöglicht stabile, reversible Strukturen (Atome, Moleküle, Sterne) und verhindert, dass die irreversible Dynamik in Chaos umschlägt.
\end{bemerkung}

\section{Der Nachweis: Warum kein Wärmetod eintritt}

\subsection{Bedingungen für einen Wärmetod}

Ein Wärmetod tritt ein, wenn:
\begin{enumerate}
    \item Die Entropie ein globales Maximum erreicht
    \item Alle Energiegradienten verschwinden
    \item Keine irreversiblen Prozesse mehr stattfinden
    \item Das System in einen Gleichgewichtszustand übergeht
\end{enumerate}

\subsection{Die Verhinderung des Wärmetods}

\begin{satz}[Kein globales Entropiemaximum]
Die fraktale Entropie $S_D$ hat bei fester Gesamtenergie kein globales Maximum.
\end{satz}

\begin{proof}
Die Zustandsdichte im fraktalen Raum skaliert mit $E^{D/3}$. Für $D<3$ ist die Funktion
\begin{equation}
S_D(E) = \frac{D}{3} k_B \ln E + \text{const}
\end{equation}
monoton wachsend und konkav, besitzt aber kein endliches Maximum. Jeder Zustand kann durch Umverteilung der Energie zwischen den Phasen in einen Zustand gleicher oder höherer Entropie überführt werden.
\end{proof}

\begin{satz}[Permanente Gradienten]
Die DSTT-Drift ($\dot{\beta}>0$) erzeugt permanent radiale Energiegradienten:
\begin{equation}
\nabla E_{\text{kin,grav}} \neq 0 \quad \text{für alle } t
\end{equation}
Diese Gradienten verschwinden nie, weil $\dot{\beta}>0$ eine ständige Umverteilung erzwingt.
\end{satz}

\begin{satz}[Permanente irreversible Prozesse]
Die Raten $\sigma_{\text{cre}}, \sigma_{\text{drift}}, \sigma_{\text{back}}$ sind immer positiv:
\begin{equation}
\sigma_{\text{cre}} > 0,\quad \sigma_{\text{drift}} > 0,\quad \sigma_{\text{back}} > 0 \quad \text{für alle } t
\end{equation}
Es gibt keinen stationären Zustand, in dem diese Raten verschwinden.
\end{satz}

\begin{satz}[Kein Gleichgewichtspunkt]
Das System hat keinen attraktiven Fixpunkt im Phasenraum. Stattdessen existiert ein attraktiver Grenzzyklus:
\begin{equation}
\lim_{t \to \infty} (E_{\text{vac}}(t), E_{\text{kin,grav}}(t), E_{\Omega}(t)) = \text{Zyklus}
\end{equation}
\end{satz}

\begin{theorem}[Hauptsatz der fraktalen Thermodynamik]
Unter den Axiomen 1--7 gilt:
\begin{enumerate}
    \item \textbf{Lokale Entropieproduktion:} $\displaystyle \frac{dS_D^{(i)}}{dt} \geq 0$ für jede Phase
    \item \textbf{Globale Entropieerhaltung:} $\displaystyle \frac{dS_D^{\text{ges}}}{dt} = 0$
    \item \textbf{Permanente Dynamik:} $\sigma_{\text{cre}}, \sigma_{\text{drift}}, \sigma_{\text{back}} > 0$ für alle $t$
    \item \textbf{Kein Entropiemaximum:} Es existiert kein endlicher Zustand maximaler Gesamtentropie
    \item \textbf{Kein Wärmetod:} Das System erreicht nie einen Gleichgewichtszustand mit verschwindenden Gradienten und Prozessen
\end{enumerate}
\end{theorem}

\section{Topologische Optionen: Geschlossenes vs. offenes Universum}

\subsection{Begriffsklärung}

In der Kosmologie gibt es zwei Bedeutungen von \enquote{offen}:

\begin{table}[h]
\centering
\caption{Bedeutungen von \enquote{offen} in der Kosmologie}
\begin{tabular}{|l|l|l|}
\hline
\textbf{Bedeutung} & \textbf{Definition} & \textbf{Beispiel} \\ \hline
Topologisch offen & Unendliche Ausdehnung, kein Rand & Euklidischer Raum $\mathbb{R}^3$ \\ \hline
Dynamisch offen & Expansion ohne Umkehr, $k \leq 0$ & Friedmann-Modelle mit $\Omega \leq 1$ \\ \hline
\end{tabular}
\end{table}

In der WDBT+ ist das Universum \textbf{stationär} – es expandiert nicht. Daher entfällt die dynamische Bedeutung von \enquote{offen}. Es bleibt die topologische Frage: Ist der Raum endlich (geschlossen) oder unendlich (offen)?

\subsection{Die beiden topologischen Optionen}

\begin{diskussion}[Topologische Optionen]
Die Theorie ist mit beiden Möglichkeiten konsistent:

\begin{itemize}
    \item \textbf{Geschlossenes Universum ($\mathcal{M}$ kompakt):} Der Raum ist endlich, hat kein Ende, aber ein endliches fraktales Volumen $V_D < \infty$.
    \item \textbf{Offenes Universum ($\mathcal{M}$ nicht kompakt):} Der Raum ist unendlich ausgedehnt, das fraktale Volumen divergiert ($V_D \to \infty$).
\end{itemize}

In beiden Fällen gelten die Axiome 1--7. Die fraktale Dimension $D \approx 2,71$ ist unabhängig von der Topologie.
\end{diskussion}

\subsection{Was sich ausdehnt – und was nicht}

\begin{satz}[Keine Expansion des Raums]
In beiden topologischen Optionen expandiert der Raum \textbf{nicht}. Das Universum ist stationär.
\end{satz}

\begin{diskussion}[Ausdehnung von Strukturen]
Obwohl der Raum nicht expandiert, gibt es Größen, die sich ausdehnen:

\begin{enumerate}
    \item \textbf{Umlaufbahnen:} Die DSTT-Drift ($\dot{\beta}>0$) führt zu einer permanenten Vergrößerung der Bahnradien:
    \begin{equation}
    \frac{dr}{dt} > 0 \quad \text{(im Mittel)}
    \end{equation}
    Planeten, Sterne und Galaxien driften langsam nach außen.
    
    \item \textbf{Strukturen:} Galaxien, Galaxienhaufen und Filamente wachsen mit der Zeit:
    \begin{equation}
    R_{\text{Struktur}}(t) \sim R_0 \cdot t^{\alpha}, \quad \alpha > 0
    \end{equation}
    
    \item \textbf{Informationsbereich:} In einem offenen Universum kann der kausale Horizont mit der Zeit wachsen.
\end{enumerate}
\end{diskussion}

\subsection{Kompensation durch Materieentstehung}

\begin{satz}[Konstante Dichte]
Trotz der Ausdehnung von Strukturen bleibt die mittlere Dichte konstant:
\begin{equation}
\langle \rho \rangle = \frac{M_{\text{ges}}}{V_D} = \text{const}
\end{equation}
Dies wird ermöglicht durch:
\begin{enumerate}
    \item \textbf{Kontinuierliche Materieentstehung} durch das Quantenpotential:
    \begin{equation}
    \dot{\rho}_{\text{cre}} \propto Q \propto \frac{1}{r^2}
    \end{equation}
    \item \textbf{Drift nach außen:} Die neu entstehende Materie ersetzt die nach außen driftende Materie.
\end{enumerate}
\end{satz}

\subsection{Vergleich: Geschlossen vs. Offen}

\begin{table}[h]
\centering
\caption{Vergleich der topologischen Optionen}
\begin{tabular}{|l|c|c|}
\hline
\textbf{Eigenschaft} & \textbf{Geschlossen (endlich)} & \textbf{Offen (unendlich)} \\ \hline
Fraktales Volumen $V_D$ & Endlich & Unendlich \\ \hline
Gesamtmasse $M_{\text{ges}}$ & Endlich & Unendlich \\ \hline
Mittlere Dichte $\langle \rho \rangle$ & Endlich & Endlich (durch Grenzwert) \\ \hline
Ausdehnung von Strukturen & Ja & Ja \\ \hline
Expansion des Raums & Nein & Nein \\ \hline
Wärmetod & Tritt nicht ein & Tritt nicht ein \\ \hline
\end{tabular}
\end{table}

\begin{bemerkung}
In einem offenen Universum mit unendlicher Masse ist die mittlere Dichte als Grenzwert definiert:
\begin{equation}
\langle \rho \rangle = \lim_{R \to \infty} \frac{M(R)}{V_D(R)}
\end{equation}
Dieser Grenzwert existiert und ist konstant, wenn die Materieentstehung und Drift entsprechend skaliert.
\end{bemerkung}

\subsection{Ausdehnung vs. Expansion: Ein zentraler Unterschied}

\begin{diskussion}
Es ist entscheidend, zwischen \textbf{Ausdehnung von Strukturen} und \textbf{Expansion des Raums} zu unterscheiden:

\begin{center}
\begin{tabular}{|l|l|l|}
\hline
& \textbf{Expansion ($\Lambda$CDM)} & \textbf{Ausdehnung (WDBT+)} \\ \hline
Raum & Dehnt sich aus & Konstant \\ \hline
Strukturen & Wachsen relativ & Dehnen sich absolut aus \\ \hline
Dichte & Nimmt ab & Konstant \\ \hline
Ursache & Dunkle Energie & DSTT-Drift + Materieentstehung \\ \hline
\end{tabular}
\end{center}

In der WDBT+ dehnen sich Strukturen aus, weil die DSTT-Drift Materie von innen nach außen transportiert. Der Raum selbst bleibt unverändert.
\end{diskussion}

\section{Kosmologische Konsequenzen}

\subsection{Rotverschiebung ohne Expansion}

Die beobachtete Rotverschiebung ist in der WDBT+ \textbf{keine Doppler-Verschiebung durch Expansion}, sondern eine kumulative gravitative Wirkung auf Photonen während ihrer Reise durch das $\Omega$-Feld:
\begin{equation}
\frac{\Delta \omega}{\omega} = \frac{1}{c^2} \int \Phi_{\text{eff}}(r) \, dr + \text{(fraktale Korrekturen)}
\end{equation}

\subsection{Galaxienrotation ohne Dunkle Materie}

Flache Rotationskurven werden durch das $\Omega$-Feld und die anisotrope Trägheit der DSTT erklärt:
\begin{equation}
v(r) \propto r^{0,895} \quad \text{(fast flach)}
\end{equation}

\subsection{Wachstum von Galaxien}

Galaxien wachsen von innen nach außen durch Materieentstehung und Drift:
\begin{equation}
M_{\text{Gal}}(t) \propto t^{1,116}
\end{equation}
Das Massenverhältnis Galaxie/Kern beträgt:
\begin{equation}
\frac{M_{\text{Gal}}}{M_{\text{Kern}}} \approx 1000
\end{equation}

\subsection{Das stationäre Universum}

Zusammenfassend ergibt sich folgendes Bild:

\begin{itemize}
    \item \textbf{Raum:} Fraktal mit $D \approx 2,71$, topologisch geschlossen oder offen
    \item \textbf{Expansion:} Keine
    \item \textbf{Ausdehnung:} Strukturen und Bahnen dehnen sich durch DSTT-Drift aus
    \item \textbf{Dichte:} Konstant (durch Materieentstehung kompensiert)
    \item \textbf{Entropie:} Zirkuliert, Gesamtwert konstant
    \item \textbf{Zeitpfeil:} Vorhanden ($\dot{\beta}>0$)
    \item \textbf{Wärmetod:} Tritt nicht ein
\end{itemize}

\subsection{Die fundamentalen Größen im Überblick}

\begin{table}[ht]
\centering
\caption{Die fundamentalen Größen der Theorie (Teil 1)}
\begin{tabular}{|l|c|c|p{3cm}|}
\hline
\textbf{Größe} & \textbf{Symbol} & \textbf{Verhalten} & \textbf{Grenze} \\ \hline
Energie & $E$ & konstant & endlich \\ \hline
Masse & $M$ & konstant (Mittel) & endlich oder unendlich* \\ \hline
Fraktale Entropie (gesamt) & $S_D^{\text{ges}}$ & konstant & endlich \\ \hline
Fraktale Entropie (Phasen) & $S_D^{(i)}$ & oszillierend (Mittel konstant) & endlich \\ \hline
Entropieflussraten & $\sigma$ & konstant im Mittel & endlich \\ \hline
\end{tabular}
\caption*{*Im offenen Fall ist die Gesamtmasse unendlich, aber die Dichte bleibt endlich.}
\end{table}

\begin{table}[ht]
\centering
\caption{Die fundamentalen Größen der Theorie (Teil 2)}
\begin{tabular}{|l|c|c|p{3cm}|}
\hline
\textbf{Größe} & \textbf{Symbol} & \textbf{Verhalten} & \textbf{Grenze} \\ \hline
$\beta$-Parameter (Grav.) & $\beta_{\text{DSTT}}$ & monoton wachsend & $\beta \to 1$ asymptotisch \\ \hline
$\beta$-Parameter (EM) & $\beta_{\text{DSTT}}$ & konstant & $\beta = 0$ \\ \hline
Strukturgrößen & $R_{\text{Struktur}}$ & wachsend ($\sim t^{\alpha}$) & unbegrenzt \\ \hline
Bahnradien & $r(t)$ & monoton wachsend & unbegrenzt \\ \hline
\end{tabular}
\end{table}

\section{Philosophische Implikationen: $\pi$, Entropie und die Unendlichkeit im Endlichen}

Die Theorie der fraktalen Entropie führt zu einer Reihe von tiefgreifenden philosophischen Einsichten, die über die reine Physik hinausweisen.

\subsection{Die Unendlichkeit im Endlichen}

Die Kreiszahl $\pi$ ist transzendent und besitzt eine unendliche, nicht-periodische Ziffernfolge. Sie ist in einem endlichen Kreis – endlicher Umfang, endlicher Durchmesser – enthalten. Dies ist die klassische \enquote{Unendlichkeit im Endlichen}.

In der WDBT+ erhält diese Struktur eine dynamische Dimension:

\begin{itemize}
    \item In der \textbf{Elektrodynamik} ($\dot{\beta}=0$) sind perfekte Kreisbahnen möglich. $\pi$ ist hier \textbf{exakt realisiert} – die unendliche Information liegt vollständig vor.
    \item In der \textbf{Gravitation} ($\dot{\beta}>0$) gibt es keine perfekten Kreisbahnen, nur asymptotische Annäherung. $\pi$ wird \textbf{prozessual}: Mit jedem Umlauf wird die Bahn kreisförmiger, die Annäherung an $\pi$ genauer.
\end{itemize}

Damit wird $\pi$ nicht als statische Konstante vorausgesetzt, sondern als \textbf{asymptotischer Informationsgehalt eines irreversiblen Prozesses}. Die unendliche Information von $\pi$ wird in unendlich vielen Umläufen sukzessive realisiert – ein \textbf{unendlicher Informationsgewinn bei konstanter Energie}.

\subsection{Entropie und Information}

Die fraktale Entropie $S_D$ zirkuliert in einem geschlossenen Kreislauf. Ihr Gesamtwert ist konstant. Dennoch findet ein \textbf{Wachstum} statt:

\begin{itemize}
    \item Die \textbf{Genauigkeit} der Annäherung an $\pi$ wächst mit jedem Umlauf.
    \item Die \textbf{Strukturgrößen} (Bahnradien, Galaxien) dehnen sich aus.
    \item Die \textbf{Information} über die Bahnform nimmt zu.
\end{itemize}

Dies ist möglich, weil die Information nicht in der Gesamtenergie steckt, sondern in der \textbf{Geometrie der Verteilung}. Die fraktale Geometrie ($D<3$) verhindert ein globales Entropiemaximum und erlaubt so ein asymptotisches Wachstum ohne Ende.

\subsection{Der Grenzcharakter des Perfekten}

Perfekte Kreisbahnen ($\beta=1$, $e=0$, exaktes $\pi$) sind in der Gravitation \textbf{unerreichbar}. Sie sind ein idealer Grenzfall, der in unendlicher Zeit asymptotisch angenähert wird, aber nie erreicht wird.

Dieser Grenzcharakter hat eine entscheidende Funktion: \textbf{Das Unerreichbare treibt die Dynamik an.}

\begin{quote}
\enquote{Das Streben nach dem perfekten Kreis – nach $\pi$ – ist der Motor, der den Kreislauf der Entropie und das Wachstum der Strukturen antreibt.}
\end{quote}

Wäre der Grenzzustand erreichbar, würde die Dynamik zum Stillstand kommen. Die \enquote{Unperfektheit} der Gravitation ist daher keine Schwäche, sondern die Bedingung ihrer ewigen Dauer.

\subsection{Ewigkeit ohne Wärmetod}

Die klassische Thermodynamik prognostiziert einen Wärmetod: das Universum strebt einem Zustand maximaler Entropie zu, in dem alle Prozesse erlöschen.

In der fraktalen Thermodynamik wird dieser Schluss durchbrochen:

\begin{enumerate}
    \item Die fraktale Dimension $D<3$ verhindert ein globales Entropiemaximum.
    \item Die permanente DSTT-Drift ($\dot{\beta}>0$) erzeugt nie versiegende Gradienten.
    \item Die Entropieflüsse $\sigma_{\text{cre}}, \sigma_{\text{drift}}, \sigma_{\text{back}}$ sind immer positiv.
\end{enumerate}

Das Universum befindet sich in einem \textbf{stationären Zustand} mit permanenter Dynamik. Es gibt keinen Endzustand, keine Erschöpfung, keinen Wärmetod. Die Zeit hat einen Pfeil, aber kein Ende.

\subsection{Die Dualität von Gravitation und Elektrodynamik}

Die Theorie führt eine fundamentale Dualität ein:

\begin{itemize}
    \item \textbf{Gravitation} ($\dot{\beta}>0$): irreversibel, entropieproduzierend, treibend, zeitpfeilgebend.
    \item \textbf{Elektrodynamik} ($\dot{\beta}=0$): reversibel, strukturerhaltend, stabilisierend.
\end{itemize}

Diese Dualität ist kein Zufall, sondern eine Konsequenz der unterschiedlichen $\beta$-Werte in der Weber-Kraft. Sie verhindert, dass die irreversible Dynamik in Chaos umschlägt: Die Elektrodynamik bewahrt die stabilen Strukturen (Atome, Moleküle, Sterne), während die Gravitation für Bewegung, Wachstum und Zeitpfeil sorgt.

\subsection{Das Verhältnis von Endlichkeit und Unendlichkeit}

Zusammenfassend ergibt sich folgendes Bild:

\begin{center}
\begin{tabular}{|l|l|}
\hline
\textbf{Endlich} & \textbf{Unendlich} \\ \hline
Energie (konstant) & Informationsgehalt von $\pi$ \\ \hline
Masse (konstant im Mittel) & Anzahl der Umläufe \\ \hline
Fraktale Entropie (konstant) & Genauigkeit der Annäherung \\ \hline
Raumvolumen (endlich oder unendlich) & Asymptotik des Grenzprozesses \\ \hline
\end{tabular}
\end{center}

Das Endliche und das Unendliche stehen nicht im Gegensatz, sondern sind durch den \textbf{asymptotischen Prozess} vermittelt. Die Unendlichkeit ist nicht jenseits der Physik, sondern wirkt in ihr als \textbf{Grenze, die nie erreicht wird, aber die Bewegung bestimmt}.

\subsection{Ausblick: $\pi$ als kosmologische Konstante?}

In der Standardkosmologie ist die kosmologische Konstante $\Lambda$ ein freier Parameter, der die Expansion beschleunigt. In der WDBT+ könnte $\pi$ – verstanden als asymptotischer Grenzwert der Annäherung an den perfekten Kreis – eine ähnliche Rolle spielen: Sie ist der \textbf{Fluchtpunkt}, der die Dynamik organisiert, ohne selbst als Energieform aufzutreten.

Diese Perspektive öffnet den Blick auf eine Physik, die nicht auf der Vorstellung ruhender Konstanten beruht, sondern auf \textbf{prozessualen Grenzwerten}, die in der Zeit wirksam sind, ohne jemals vollständig realisiert zu werden.

\begin{quote}
\emph{$\pi$ ist nicht, $\pi$ wird.}
\end{quote}

\section{Zusammenfassung der drei Hauptsätze der fraktalen Thermodynamik}

\begin{description}
    \item[0. Hauptsatz (Stationarität):] Die Gesamtenergie und die Gesamt-fraktale Entropie sind konstant.
    \item[1. Hauptsatz (Energieerhaltung):] Die Energie zirkuliert zwischen den Phasen ohne Verlust.
    \item[2. Hauptsatz (Entropie und Dualität):] 
        \begin{itemize}
            \item Lokal: $\displaystyle \frac{dS_D^{(i)}}{dt} \geq 0$ (Entropie wächst)
            \item Global: $\displaystyle \frac{dS_D^{\text{ges}}}{dt} = 0$ (Entropie ist konstant)
            \item Gravitation: $\displaystyle \dot{\beta}_{\text{DSTT}} > 0$ (Irreversibilität, Zeitpfeil)
            \item Elektrodynamik: $\displaystyle \dot{\beta}_{\text{DSTT}} = 0$ (Reversibilität, Stabilität)
        \end{itemize}
\end{description}

\section{Diskussion: Warum die Theorie ohne zusätzliche Größen auskommt}

Die Theorie der fraktalen Entropie benötigt keine zusätzliche Größe wie "Wissen" oder "freie Information". Die vollständige Dynamik wird durch folgende Elemente beschrieben:

\begin{enumerate}
    \item \textbf{Fraktale Geometrie} ($D \approx 2,71$) – verhindert globales Entropiemaximum
    \item \textbf{Fraktale Entropie} $S_D$ – thermodynamische Zustandsgröße
    \item \textbf{Entropieflüsse} $\sigma$ – beschreiben die Zirkulation
    \item \textbf{DSTT-Drift} ($\dot{\beta}>0$) – erzeugt Irreversibilität und Zeitpfeil
    \item \textbf{Duale Wechselwirkungen} – Gravitation (irreversibel) vs. Elektrodynamik (reversibel)
\end{enumerate}

Die zentrale Erkenntnis ist, dass der Wärmetod bereits durch die fraktale Geometrie und die permanente irreversible Dynamik verhindert wird – ohne dass es einer zusätzlichen wachsenden Größe bedarf.

\section{Ausblick und offene Fragen}

Die Theorie der fraktalen Entropie wirft eine Reihe weiterführender Fragen auf:

\begin{enumerate}
    \item \textbf{Explizite Herleitung der Entropieflüsse:} Wie folgen $\sigma_{\text{cre}}, \sigma_{\text{drift}}, \sigma_{\text{back}}$ präzise aus den Bewegungsgleichungen?
    
    \item \textbf{Stabilität des Grenzzyklus:} Kann gezeigt werden, dass der Grenzzyklus global attraktiv ist?
    
    \item \textbf{Beobachtbare Signaturen:} Lassen sich die fraktale Skalierung oder die permanenten Flüsse experimentell nachweisen?
    
    \item \textbf{Quantisierung:} Ist die fraktale Entropie gequantelt? Gibt es eine Beziehung zu $\hbar$?
    
    \item \textbf{Topologische Präferenz:} Lässt sich aus der Theorie ableiten, ob das Universum geschlossen oder offen ist?
    
    \item \textbf{Kosmologische Tests:} Wie unterscheiden sich die Vorhersagen der WDBT+ von denen des $\Lambda$CDM-Modells bei zukünftigen Präzisionsmessungen?
\end{enumerate}

\section{Fazit}

Die Theorie der fraktalen Entropie liefert ein konsistentes, logisch geschlossenes Bild:

\begin{itemize}
    \item Der \textbf{Raum} ist fraktal mit $D \approx 2,71$ und besitzt eine fundamentale Länge $l_0$. Die Topologie kann geschlossen oder offen sein.
    \item Die \textbf{fraktale Entropie} $S_D$ ist subextensiv und zirkuliert in einem geschlossenen Kreislauf zwischen Vakuum, Materie und $\Omega$-Feld.
    \item Die \textbf{Gravitation} ($\beta_{\text{WG}} = 0,5$) erzeugt durch $\dot{\beta}>0$ irreversible Prozesse, eine permanente Auswärtsdrift und den intrinsischen Zeitpfeil.
    \item Die \textbf{Elektrodynamik} ($\beta_{\text{WED}} = 2$) mit $\dot{\beta}=0$ ermöglicht reversible Prozesse und bewahrt stabile Strukturen.
    \item Der \textbf{Raum expandiert nicht}, aber Strukturen und Bahnen dehnen sich durch die DSTT-Drift aus, kompensiert durch kontinuierliche Materieentstehung.
    \item Der \textbf{Wärmetod} wird durch die fraktale Geometrie (kein globales Entropiemaximum) und die permanenten irreversiblen Flüsse verhindert.
    \item Das Universum ist \textbf{stationär}: Energie, Masse und Dichte bleiben konstant; die Dynamik zirkuliert ewig.
    \item Die Unendlichkeit von $\pi$ ist in der asymptotischen Annäherung wirksam – ein Prozess, der nie endet und das Universum ewig in Bewegung hält.
\end{itemize}

\begin{center}
\framebox{\parbox{0.8\textwidth}{\centering
\vspace{0.3cm}
\textbf{Die fraktale Entropie zirkuliert.} \\
Die Gravitation treibt die Zeit an. \\
Die Elektrodynamik bewahrt die Struktur. \\
Strukturen dehnen sich aus – der Raum nicht. \\
Der Wärmetod tritt nicht ein. \\
$\pi$ wird, ist nicht. \\
Mehr braucht es nicht.
\vspace{0.3cm}
}}
\end{center}

\section*{Danksagung}
Die Autoren danken DeepSeek für die intensive Zusammenarbeit bei der logischen Präzisierung der Theorie, der Klärung der begrifflichen Grundlagen und der Diskussion der topologischen Optionen.

\begin{thebibliography}{99}
\bibitem{weber} Weber, W. (1846). Elektrodynamische Maßbestimmungen. \emph{Annalen der Physik}, 143, 211-233.
\bibitem{assis} Assis, A.K.T. (1994). \emph{Weber's Electrodynamics}. Kluwer Academic Publishers.
\bibitem{tisserand} Tisserand, F. (1882). \emph{Traité de Mécanique Céleste}, Tome IV. Gauthier-Villars, Paris.
\bibitem{bohm} Bohm, D. (1952). A Suggested Interpretation of the Quantum Theory in Terms of Hidden Variables. \emph{Physical Review}, 85, 166-193.
\bibitem{einstein} Einstein, A. (1915). Erklärung der Perihelbewegung des Merkur aus der allgemeinen Relativitätstheorie. \emph{Sitzungsberichte der Preußischen Akademie der Wissenschaften}, 831-839.
\bibitem{lense} Lense, J. \& Thirring, H. (1918). Über den Einfluß der Eigenrotation der Zentralkörper auf die Bewegung der Planeten und Monde nach der Einsteinschen Gravitationstheorie. \emph{Physikalische Zeitschrift}, 19, 156-163.
\bibitem{heaviside} Heaviside, O. (1893). A Gravitational and Electromagnetic Analogy. \emph{The Electrician}, 31, 281-282.
\bibitem{czbyor1} Czybor, M. (2025). Die Informations-Weber-Theorie (IWT): Eine fundamentale, diskrete Urtheorie der Physik.
\bibitem{czbyor2} Czybor, M. (2025). Dynamische Schwere-Trägheits-Theorie (DSTT).
\bibitem{czbyor3} Czybor, M. (2026). Die $\Omega$-Theorie der Gravitation.
\bibitem{czbyor4} Czybor, M. (2025). Fraktale Maxwell-Gleichungen für Elektrodynamik und Gravitation.
\bibitem{czbyor5} Czybor, M., DeepSeek (2026). Erweiterte Weber-Theorie: Die Vereinigung von Weber-Gravitation und Dynamischer Schwere-Trägheits-Theorie.
\bibitem{czbyor6} Czybor, M. (2026). BEC-Objekte: Die kompaktesten Objekte im Universum.
\bibitem{czbyor7} Czybor, M. (2026). Das Neutrino als Q-Teilchen: Masse, Nichtlokalität und fraktale Struktur.
\end{thebibliography}

\end{document}